% Part: normal-modal-logic
% Chapter: filtrations
% Section: introduction

\documentclass[../../../include/open-logic-section]{subfiles}

\begin{document}

\olfileid{nml}{fil}{int}

\olsection{పరిచయం}

ఒక తర్కం గురించి ముఖ్యమైన ప్రశ్న: అది నిర్ణయించదగినదేనా?
అంటే, ``ఈ \tetoken{సూత్రం}{!!{formula}} చెల్లుబాటవుతుందా?''
అనే ప్రశ్నకు జవాబిచ్చే కార్యసాధ్య ప్రక్రియ ఉందా?
ప్రతిజ్ఞావాక్య తర్కం నిర్ణయించదగినది: సత్య పట్టికను
నిర్మించి \tetoken{ఒక సూత్రం}{!!a{formula}} సర్వసత్యమో
కాదో పరీక్షించవచ్చు; ఇచ్చిన \tetoken{సూత్రానికి}{!!{formula}}
ఆ పట్టిక పరిమితమే. కానీ మోడల్ సూత్రం అన్ని నమూనాల్లో
సత్యమో కాదో నేరుగా పరీక్షించలేం, ఎందుకంటే నమూనాలు
అనంతంగా ఉన్నాయి. ఇచ్చిన \tetoken{సూత్రానికి}{!!{formula}} సంబంధించిన అన్ని
పరిమిత నమూనాలను లెక్కించవచ్చు; అందులో నిజంగా వచ్చే
\tetoken{ప్రతిజ్ఞావాక్య చరరాశులకు}{!!{propositional variable}s}
లోకాల ఉపసమితులను కేటాయించడమే ముఖ్యం. ప్రాప్యత సంబంధం
స్థిరంగా ఉంటే, సాధ్యమైన వేర్వేరు కేటాయింపులు $V(p)$
అనేవి~$W$ ఉపసమితులన్నీ మాత్రమే; $\card{W} = n$ అయితే
అవి $2^n$. మన \tetoken{సూత్రం}{!!{formula}}~$!A$లో $m$
\tetoken{ప్రతిజ్ఞావాక్య చరరాశులు}{!!{propositional
  variable}s} ఉంటే, $n$ లోకాలతో కూడిన వేర్వేరు నమూనాలు
$2^{nm}$. ప్రతిదానిలోనూ ప్రతి లోకంలో~$!A$ సత్యమో కాదో
దాని సత్యమూల్యాన్ని లెక్కించి పరీక్షించవచ్చు. అన్ని
సాధ్యమైన ప్రాప్యత సంబంధాలనూ పరీక్షించాలి; అయితే $n$
లోకాలపై సంబంధాలు కూడా పరిమిత సంఖ్యలోనే ఉన్నాయి
(ఖచ్చితంగా, $W \times W$ ఉపసమితుల సంఖ్య, అంటే
$2^{n^2}$).

మన ఆసక్తి \Log{K}లో కాక, ఒక నమూనా వర్గం నిర్వచించే
తర్కంలో ఉంటే (ఉదాహరణకు, స్వావర్తన, సంక్రామక నమూనాలు),
ప్రాప్యత సంబంధం సరైన రకానిదో కూడా పరీక్షించగలగాలి.
మనకు ఆసక్తిగల ఫ్రేమ్‌లను మోడల్ \tetoken{సూత్రాల}{!!{formula}s}
ద్వారా నిర్వచించగలిగితే ఇది సాధ్యం (ఉదాహరణకు,
ఫ్రేమ్‌లో \Ax{T}, \Ax{4} చెల్లుబాటును పరీక్షించడం ద్వారా).
కాబట్టి అన్ని పరిమిత ఫ్రేమ్‌లను ఒక్కొక్కటిగా తీసుకుని,
ప్రతి ఫ్రేమ్ మన వర్గంలో ఉందో లేదో పరీక్షించి, దానిపై
సాధ్యమైన నమూనాలన్నిటినీ జాబితా చేసి, ప్రతి నమూనాలో
$!A$ సత్యమో కాదో పరీక్షించాలనే ఆలోచన. కాకపోతే ఆపవచ్చు:
$!A$ మనకు ఆసక్తిగల నమూనాల వర్గంలో చెల్లుబాటు కాదు.

ఈ ఆలోచనలో సమస్య ఉంది: వెతకడాన్ని ఎప్పుడు ఆపవచ్చో,
అసలు ఆపవచ్చో లేదో తెలియదు. \tetoken{సూత్రానికి}{!!{formula}}
పరిమిత ప్రతినమూనా ఉంటే మన ప్రక్రియ దాన్ని కనుగొంటుంది.
కానీ పరిమిత ప్రతినమూనా లేకపోతే సమాధానం రాదు.
\tetoken{సూత్రం}{!!{formula}} చెల్లుబాటవవచ్చు
(ప్రతినమూనాలే ఉండవు), లేదా మనం ఎన్నడూ చూడని అనంత
ప్రతినమూనా మాత్రమే ఉండవచ్చు. ప్రతినమూనా ఉన్న ప్రతి
\tetoken{సూత్రానికి}{!!{formula}} పరిమిత ప్రతినమూనా
ఉంటుందని చూపగలిగితే ఈ సమస్యను అధిగమించవచ్చు.
అప్పుడు ఆ తర్కానికి \emph{పరిమిత నమూనా ధర్మం} ఉందంటాం.

అయితే ఒక తర్కానికి పరిమిత నమూనా ధర్మం ఉందని ఎలా
చూపాలి? ఒక మార్గం,~$!A$కు చెందిన అనంత (ప్రతి)నమూనాను
పరిమిత నమూనాగా మార్చడం. అది సాధ్యమైతే, $!A$ సత్యం
కాని నమూనా ఉన్నప్పుడల్లా, దాని నుంచి వచ్చే పరిమిత
నమూనాలోనూ $!A$ సత్యం కాదు. ఆ పరిమిత నమూనా మన అన్ని
పరిమిత నమూనాల జాబితాలో వస్తుంది; చివరకు చెల్లుబాటు
కాని ప్రతి సూత్రం చెల్లుబాటు కాదని తెలుసుకుంటాం. కానీ
సూత్రం చెల్లుబాటైతే మన ప్రక్రియ ఆగదు. అదనంగా,
ప్రక్రియ ఇచ్చే పరిమిత నమూనా పరిమాణానికి ఒక గరిష్ఠ
హద్దు ఉందనీ, ఆ హద్దు \tetoken{సూత్రం}{!!{formula}}~$!A$
పైనే ఆధారపడుతుందనీ చూపగలిగితే, ప్రతినమూనాల కోసం
ఎంత పరిమాణం వరకు పరిమిత నమూనాలను పరీక్షించాలో
తెలుస్తుంది. అప్పటికి ప్రతినమూనా దొరకకపోతే ఏదీ లేదు.
అప్పుడు ఏ \tetoken{సూత్రం}{!!{formula}}~$!A$కైనా
``$!A$ చెల్లుబాటవుతుందా?'' అనే ప్రశ్నను మన ప్రక్రియ
నిజంగానే నిర్ణయిస్తుంది.

అనంత నిర్మాణాలను పరిమిత నిర్మాణాలుగా మార్చడానికి
తరచూ ఉపయోగపడే వ్యూహం, సంబంధిత అంశాల్లో ఒకేలా
ప్రవర్తించే నిర్మాణ \tetoken{మూలకాలను}{!!{element}s}
``ఒకటిగా గుర్తించడం''. $\mModel{M}$లో సంబంధిత అంశాల్లో
ఒకేలా ప్రవర్తించే అనంత లోకాలు ఉంటే, \emph{అన్ని}
లోకాలనూ అటువంటి లోకాల పరిమిత సంఖ్యలోని (కొన్నిసార్లు
అనంతమైన) ``వర్గాల''లో చేర్చవచ్చు. అంటే లోకాల సమితిని
సరైన రీతిలో విభజించాలి: ప్రతి వర్గంలో పరిమిత సంఖ్యలోనో
అనంత సంఖ్యలోనో లోకాలు ఉండవచ్చు; కానీ వర్గాల సంఖ్య
పరిమితమే.
\sourcecorrection{OLTENMLFILINT-001}{స్థిర మూలంలో ప్రతి
  విభాగంలో అనంత సంఖ్యలో లోకాలు ఉంటాయని ఉంది. అనంత
  నమూనాను పరిమిత వర్గాలుగా విభజించినా కొన్ని వర్గాల్లో
  పరిమిత సంఖ్యలో, ఒక్క లోకం మాత్రమే కూడా ఉండవచ్చు.
  కావలసినది వర్గాల సంఖ్య పరిమితంగా ఉండటమే.}
ఆ విభాగాలే లోకాలుగా ఉండే కొత్త నమూనా~$\mModel{M^*}$ను
నిర్వచిద్దాం. పాత నమూనాలో పరిమిత సంఖ్యలోని విభాగాలు
కొత్త నమూనాలో పరిమిత సంఖ్యలోని లోకాలను, అంటే పరిమిత
నమూనాను ఇస్తాయి. లోకం~$w$ ఉన్న విభాగాన్ని $[w]$
అందాం. పాత నమూనా~$!A$కు ప్రతినమూనా అయితే కొత్తదీ
అలాగే ఉండాలంటే $\mSat{M}{!A}[w]$ అయితే, అప్పుడు మాత్రమే
$\mSat{M^*}{!A}[{[w]}]$ కావాలి. ఇందుకు విభజనతోపాటు
నమూనా~$\mModel{M^*}$ ప్రాప్యత సంబంధాన్నీ తగిన విధంగా
నిర్వచించాలి.

ఇది ఎలా పనిచేస్తుందో చూడడానికి, మొదట ప్రతి లోకం
నుంచి ప్రతి లోకం ప్రాప్యమయ్యే సరళ సందర్భాన్ని
ఊహిద్దాం. $\mSat{M}{\Box !B}[w]$ అయితే, అప్పుడు మాత్రమే
ప్రతి $v \in W$కు $\mSat{M}{!B}[v]$; $\mModel{M^*}$కూ
అదే, అయితే $[w]$, $[v]$లను తీసుకుంటాం.
\sourcecorrection{OLTENMLFILINT-002}{స్థిర మూలం ఇక్కడ
  ``ప్రాప్యత సంబంధం లేని'' సందర్భం అంటూనే Box సత్యానికి
  ఏదో ఒక లోకంలో మళ్లీ Box సత్యాన్ని షరతుగా ఇచ్చింది.
  తరువాతి ఆగమన వాదనకు తగిన సరళ సందర్భం సార్వత్రిక
  ప్రాప్యత; అందులో అన్ని లోకాల్లో అంతర్గత సూత్రం
  సత్యమైతేనే Box సత్యం. దీనిని ఇక్కడ స్పష్టం చేశాం.}
మొదటి ప్రయత్నంగా, $!A$లో వచ్చే
\tetoken{ప్రతిజ్ఞావాక్య చరరాశుల}{!!{propositional variable}s}
విషయంలో రెండు లోకాలు $u$, $v$ ఏకీభవిస్తే (ఒకే
విభాగంలో ఉంటే) అవి తుల్యమని అనుకుందాం. అంటే,
$\mModel{M}$లో
$\mSat{M}{p}[u]$ అయితే, అప్పుడు మాత్రమే
$\mSat{M}{p}[v]$.
\sourcecorrection{OLTENMLFILINT-003}{స్థిర మూలం మొదట
  నమూనాలోని అన్ని ప్రతిజ్ఞావాక్య చరరాశులపై ఏకీభావాన్ని
  కోరుతుంది. చరరాశులు అనంతంగా ఉంటే అది పరిమిత
  విభాగాలను ఇవ్వకపోవచ్చు. ఇచ్చిన $!A$లో వచ్చే పరిమిత
  చరరాశులకు మాత్రమే ఈ ప్రారంభ పరీక్షను పరిమితం చేసి,
  తరువాత అన్ని ఉపసూత్రాలపై పరీక్షకు విస్తరించాం.}
$V^*(p) = \Setabs{[w]}{\mSat{M}{p}[w]}$ అని ఉంచుదాం.
$\mSat{M}{!A}[w]$ అయితే, అప్పుడు మాత్రమే
$\mSat{M^*}{!A}[{[w]}]$ అని చూపడమే లక్ష్యం. సహజంగా
ఆగమనంతో నిరూపిస్తాం. ఆధార సందర్భం $!A \equiv p$.
మొదట $\mSat{M}{p}[w]$ అనుకుందాం. నిర్వచనం వల్ల
$[w] \in V^*(p)$; కనుక $\mSat{M^*}{p}[{[w]}]$.
\sourcecorrection{OLTENMLFILINT-004}{స్థిర మూలంలోని
  ఈ మొదటి మూలకత్వ వాక్యంలో $V^*$ అని మాత్రమే ఉంది;
  పరమాణువు $p$కు సంబంధించిన విలువ సమితి $V^*(p)$
  కావాలి. తరువాతి వాక్యంలో మూలమే $V^*(p)$ను వాడింది.}
ఇప్పుడు $\mSat{M^*}{p}[{[w]}]$ అనుకుందాం. అంటే
$[w] \in V^*(p)$; కాబట్టి $w$కు తుల్యమైన ఏదో $v$కు
$\mSat{M}{p}[v]$. ``$w$, $v$కు తుల్యం'' అంటే
``$w$, $v$ సంబంధిత
\tetoken{ప్రతిజ్ఞావాక్య చరరాశుల}{!!{propositional variable}s}న్నింటినీ
ఒకేలా సత్యం చేస్తాయి''; కనుక $\mSat{M}{p}[w]$.
ఇప్పుడు ఆగమన దశలో, ఉదాహరణకు, $!A \ident \lnot !B$
అనుకుందాం. అప్పుడు $\mSat{M}{\lnot !B}[w]$ అయితే,
అప్పుడు మాత్రమే $\mSat/{M}{!B}[w]$; ఆగమన పరికల్పన
ప్రకారం, అప్పుడు మాత్రమే $\mSat/{M^*}{!B}[{[w]}]$;
అప్పుడు మాత్రమే $\mSat{M^*}{\lnot !B}[{[w]}]$.
ఇతర మోడల్-యేతర సంయోజకాలకూ ఇదే విధంగా పనిచేస్తుంది.
$\Box$కు కూడా పనిచేస్తుంది: $\mSat{M^*}{\Box
  !B}[{[w]}]$ అనుకుందాం. అంటే, ప్రతి $[u]$కు
$\mSat{M^*}{!B}[{[u]}]$. ఆగమన పరికల్పన వల్ల ప్రతి
$u$కు $\mSat{M}{!B}[u]$. అందువల్ల
$\mSat{M}{\Box !B}[w]$.

సాధారణ సందర్భంలో $\mModel{M^*}$ ప్రాప్యత సంబంధాన్ని
కూడా నిర్వచించాలి; అప్పుడు వాదన మరింత క్లిష్టమవుతుంది.
దాని ప్రాప్యత సంబంధం~$R^*$ పైన చెప్పిన ఆగమన
నిరూపణను సాధ్యం చేసే అవసరమైన షరతులను నెరవేర్చితే,
నమూనా~$\mModel{M^*}$ను \emph{వడపోత} అందాం. అప్పుడు
ఏ వడపోత~$\mModel{M^*}$లోనైనా $[w]$ వద్ద $!A$
సత్యమైతే, అప్పుడు మాత్రమే $\mModel{M}$లో~$w$ వద్ద
$!A$ సత్యం. అయితే ఇప్పుడు వడపోతలు \emph{ఉన్నాయని}
కూడా చూపాలి: అంటే, అవసరమైన షరతులను నెరవేర్చేలా~$R^*$ను
నిర్వచించగలమని చూపాలి. అందుకు లోకాలు $u$, $v$లను
కేవలం పైన పేర్కొన్న
\tetoken{ప్రతిజ్ఞావాక్య చరరాశుల}{!!{propositional variable}s}
విషయంలోనే కాక,~$!A$ యొక్క అన్ని
\tetoken{ఉపసూత్రాల}{!!{formula}s} విషయంలోనూ
ఏకీభవించినప్పుడే తుల్యంగా తీసుకోవాలి. $!A$కు పరిమిత
సంఖ్యలోనే \tetoken{ఉపసూత్రాలు}{!!{formula}s} ఉన్నందున,
వడపోత పరిమితమని ఇది ఇప్పటికీ హామీ ఇస్తుంది. వడపోతను
నిర్వచించడానికి ఒక్క మార్గమే లేదు; దాని ప్రాప్యత
సంబంధానికి కావలసిన ధర్మాలు (ఉదాహరణకు, స్వావర్తన,
సంక్రామక మొదలైనవి) ఉండేలా $R^*$ నిర్వచనంలో
సృజనాత్మకంగా వ్యవహరించాలి.


\end{document}
